% Mathématiques 4e — cours V3
% Arithmétique et nombres premiers

\section{Objectifs}

\begin{itemize}
\item Utiliser les notions de multiple et de diviseur.
\item Reconnaître un nombre premier.
\item Connaître ou retrouver les nombres premiers inférieurs ou égaux à $100$.
\item Décomposer un entier en produit de facteurs premiers.
\item Utiliser cette décomposition pour simplifier une fraction.
\item Reconnaître et produire des fractions égales.
\item Obtenir une fraction irréductible.
\item Résoudre des problèmes de divisibilité, de partage ou de périodicité.
\end{itemize}

\section{1. Multiples et diviseurs}

Pour des entiers $a$ et $b$ avec $b\neq0$, dire que $b$ divise $a$ signifie qu'il existe un entier $k$ tel que :
\[
a=bk.
\]

Par exemple :
\[
84=7\times12.
\]

Donc $7$ et $12$ sont des diviseurs de $84$.

\section{2. Critères de divisibilité}

Les critères usuels restent utiles.

Un entier est divisible :
\begin{itemize}
\item par $2$ si son chiffre des unités est pair ;
\item par $3$ si la somme de ses chiffres est divisible par $3$ ;
\item par $5$ s'il se termine par $0$ ou $5$ ;
\item par $9$ si la somme de ses chiffres est divisible par $9$ ;
\item par $10$ s'il se termine par $0$.
\end{itemize}

\section{3. Nombre premier}

\begin{definition}
Un nombre premier est un entier positif qui possède exactement deux diviseurs positifs :
\[
1
\]
et lui-même.
\end{definition}

Ainsi :
\[
2,\ 3,\ 5,\ 7,\ 11
\]
sont premiers.

Le nombre $1$ n'est pas premier.

\section{4. Nombres premiers jusqu'à cent}

\coursefigure{/images/mathematiques/4eme/premiers-grille.svg}{Grille des entiers de 1 à 100 avec nombres premiers mis en évidence.}{Les nombres premiers sont les briques élémentaires de la décomposition multiplicative des entiers.}

La liste utile jusqu'à $100$ commence par :
\[
2,\ 3,\ 5,\ 7,\ 11,\ 13,\ 17,\ 19,\ 23,\ 29,
\]
et se termine par :
\[
71,\ 73,\ 79,\ 83,\ 89,\ 97.
\]

Il faut surtout savoir les retrouver efficacement.

\section{5. Tester un nombre}

Pour :
\[
91,
\]
on observe :
\[
91=7\times13.
\]

Donc $91$ n'est pas premier.

Pour $97$, il suffit de tester les petits nombres premiers jusqu'à environ $\sqrt{97}<10$ :
\[
2,\ 3,\ 5,\ 7.
\]

Aucun ne divise $97$.

Donc $97$ est premier.

\section{6. Décomposition}

\begin{definition}
Décomposer un entier en facteurs premiers consiste à l'écrire comme produit de nombres premiers.
\end{definition}

Pour $84$ :
\[
84=2\times42
\]
puis :
\[
42=2\times21
\]
et :
\[
21=3\times7.
\]

Donc :
\[
\boxed{84=2^2\times3\times7}.
\]

\coursefigure{/images/mathematiques/4eme/premiers-arbre.svg}{Arbre de facteurs de 84 aboutissant à 2, 2, 3 et 7.}{On poursuit la décomposition jusqu'à ce que tous les facteurs obtenus soient premiers.}

\section{7. Divisions successives}

Pour $360$ :
\[
360\div2=180,
\]
\[
180\div2=90,
\]
\[
90\div2=45,
\]
\[
45\div3=15,
\]
\[
15\div3=5,
\]
\[
5\div5=1.
\]

Donc :
\[
\boxed{360=2^3\times3^2\times5}.
\]

\section{8. Simplifier une fraction}

Considérons :
\[
\dfrac{84}{126}.
\]

On a :
\[
84=2^2\times3\times7,
\]
\[
126=2\times3^2\times7.
\]

On simplifie les facteurs communs :
\[
2\times3\times7.
\]

Donc :
\[
\dfrac{84}{126}=\dfrac23.
\]

\section{9. Fraction irréductible}

\begin{definition}
Une fraction est irréductible lorsque son numérateur et son dénominateur n'ont aucun facteur premier commun.
\end{definition}

\[
35=5\times7,
\]
\[
48=2^4\times3.
\]

Donc :
\[
\dfrac{35}{48}
\]
est irréductible.

\section{10. Fractions égales}

\[
\dfrac57=\dfrac{30}{42}
\]
car on a multiplié le numérateur et le dénominateur par :
\[
6.
\]

La décomposition permet de reconnaître les facteurs communs de deux écritures.

\section{11. Problème de lots}

On dispose de :
\[
84
\]
jetons rouges et :
\[
126
\]
jetons bleus.

On veut le plus grand nombre de lots identiques.

Les facteurs communs donnent :
\[
42
\]
lots.

Chaque lot contient :
\[
84\div42=2
\]
jetons rouges et :
\[
126\div42=3
\]
jetons bleus.

\section{12. Problème de périodicité}

Deux alarmes sonnent toutes les :
\[
18\ \text{min}
\]
et :
\[
24\ \text{min}.
\]

\[
18=2\times3^2,
\]
\[
24=2^3\times3.
\]

Le premier multiple commun utile est :
\[
2^3\times3^2=72.
\]

Elles sonnent ensemble après :
\[
\boxed{72\ \text{min}}.
\]

\section{13. Algorithme simple}

\begin{verbatim}
tester les diviseurs premiers
si un diviseur exact est trouvé
    le nombre n'est pas premier
sinon
    le nombre est premier
\end{verbatim}

Un programme permet d'automatiser des tests mais ne remplace pas la compréhension des critères.

\section{14. Recherche efficace}

Si :
\[
n=ab,
\]
au moins un des facteurs est inférieur ou égal à :
\[
\sqrt n.
\]

On peut donc limiter les tests de divisibilité à des facteurs suffisamment petits.

\section{15. Exemple 210/462}

\[
210=2\times3\times5\times7,
\]
\[
462=2\times3\times7\times11.
\]

On simplifie par :
\[
2\times3\times7=42.
\]

Donc :
\[
\dfrac{210}{462}=\dfrac5{11}.
\]

\section{16. Choisir une méthode}

Pour tester la divisibilité par $5$, un critère suffit.

Pour simplifier une fraction complexe, la décomposition en facteurs premiers peut être plus claire.

La méthode dépend donc du problème.

\section{17. Méthode générale}

\begin{methode}
\begin{enumerate}
\item tester les plus petits nombres premiers ;
\item diviser tant que le quotient est entier ;
\item poursuivre jusqu'à obtenir uniquement des facteurs premiers ;
\item écrire le produit obtenu ;
\item utiliser ces facteurs pour simplifier ou résoudre le problème.
\end{enumerate}
\end{methode}

\section{18. Erreurs fréquentes}

\begin{itemize}
\item Considérer $1$ comme premier.
\item Arrêter une décomposition alors qu'un facteur est encore composé.
\item Confondre facteur premier et chiffre premier.
\item Simplifier numérateur et dénominateur par des nombres différents.
\item Croire qu'une fraction est irréductible sans vérifier les facteurs communs.
\end{itemize}

\section{À retenir}

\begin{aretenir}
Un nombre premier possède exactement deux diviseurs positifs :
\[
\boxed{1\text{ et lui-même}}.
\]

Tout entier supérieur à $1$ peut être décomposé en produit de facteurs premiers.

Cette décomposition permet de simplifier des fractions, d'obtenir une fraction irréductible et de résoudre des problèmes de divisibilité, de partage ou de périodicité.
\end{aretenir}
